% =====================================================================
% Chương 05 — Lớp Ứng dụng (UI / LVGL v9)
% Phân tích: ui_dashboard.{c,h}, ui_dashboard_layout.c, ui_dashboard_system.c,
%            ui_dashboard_private.h, ui_dashboard_theme.h
% =====================================================================

\chapter{Lớp Ứng dụng — UI (LVGL v9)}
\label{ch:application-ui}

Lớp UI của \textbf{waveshare-screen} là giao diện LVGL v9 hiển thị 6 vùng
cảm biến, đánh giá nguy cơ tổng thể và trang System.  Component được tách theo
trách nhiệm (\textbf{R7} — mỗi file $\le$ 400 dòng):

\begin{center}
\begin{longtable}{@{}p{5.2cm}p{2.4cm}p{5.6cm}@{}}
  \toprule
  \textbf{File} & \textbf{Dòng} & \textbf{Trách nhiệm} \\
  \midrule
  \endhead
  \code{ui\_dashboard.c} & 331 & API công khai + \code{evaluate\_hazard} + mute \\
  \code{ui\_dashboard\_layout.c} & 335 & builder widget (header, sidebar, canvas) \\
  \code{ui\_dashboard\_system.c} & 162 & trang System + thông tin heap/uptime/token \\
  \code{ui\_dashboard\_private.h} & 75 & khai báo biến global widget \\
  \code{ui\_dashboard\_theme.h} & 42 & bảng màu + ánh xạ zone $\rightarrow$ màu \\
  \code{include/ui\_dashboard.h} & 64 & API public (hàm gọi từ \code{main.c}) \\
  \bottomrule
\end{longtable}
\end{center}

\section{Đánh giá nguy cơ — \code{evaluate\_hazard}}
\label{sec:eval-hazard}

\textbf{Trung tâm của UI} là hàm \code{evaluate\_hazard()} (\code{ui\_dashboard.c:134}):
\begin{enumerate}
\item \code{sensor\_model\_get\_all(readings)} — snapshot an toàn.
\item \textbf{Skip các cảm biến chưa có dữ liệu} (\code{readings[i].is\_stale},
  \code{ui\_dashboard.c:145}): nếu không skip, \code{distance\_cm = 0} mặc định
  sẽ phân loại \textbf{DANGER} và ghim vĩnh viễn banner OVERALL (chú thích
  \code{ui\_dashboard.c:141--144}).
\item Với cảm biến còn lại, \code{sensor\_model\_classify} theo zone chung R3;
  \code{worst = max(zone)}.
\item Cập nhật \code{s\_lbl\_hazard\_overall}: \textbf{OVERALL: DANGER /
  CAUTION / SAFE}; nếu DANGER và \code{s\_alarm\_muted} thì ghi thêm
  \code{(MUTED)} (\code{ui\_dashboard.c:153--160}).
\end{enumerate}

\noindent\textbf{Mute:} \code{mute\_btn\_cb} (\code{ui\_dashboard.c:69}) đảo
\code{s\_alarm\_muted}, cập nhật nút ("Mute Alarm"/"Unmute Alarm") và gọi lại
\code{evaluate\_hazard()} cùng \code{arc\_set\_zone}.  Nút chỉ làm \textbf{tắt
blink/animation} trên zone DANGER, banner vẫn hiện DANGER — không thể nhầm bị
mute với SAFE (\code{ui\_dashboard.c:156--158}).

\section{Bố cục widget — \code{ui\_dashboard\_layout.c}}
\label{sec:layout}

\textbf{Các builder:}
\begin{itemize}
\item \code{build\_header} (row 90): thanh đầu với title, \code{ESP-NOW} /
  \code{Wi-Fi} / \code{MQTT} status.
\item \code{build\_left\_sidebar} (row 146): các nút hành động; \textbf{lưu ý}
  nút \textbf{"Calibrate"} (\code{ui\_dashboard\_layout.c:177--181}) được tạo
  nhưng \textbf{không} có \code{lv\_obj\_add\_event\_cb} — hiện chỉ là
  placeholder trực quan (không gọi callback).
\item \code{build\_center\_canvas} (row 190): bố cục hình chiếc xe — hood,
  cabin, trunk + 6 cung \code{arc} cho mỗi vùng cảm biến
  (\code{arc\_set\_zone} / \code{arc\_set\_nodata}).
\item \code{build\_right\_sidebar} (row 253): \code{OVERALL} hazard, crossing
  risk, relay, buzzer state.
\end{itemize}

\noindent\textbf{Bảng màu zone} (\code{ui\_dashboard\_theme.h:15--23}):
\texttt{COLOR\_DANGER 0xFF1744}, \texttt{COLOR\_CAUTION 0xFFD600},
\texttt{COLOR\_SAFE 0x00C853}, \texttt{COLOR\_NODATA 0x64748B},
\texttt{COLOR\_ACCENT 0x38BDF8}, nền \texttt{COLOR\_BG 0x0F172A};
ánh xạ zone $\rightarrow$ màu ở \code{ui\_dashboard\_theme.h:39--41}.

\section{Trang System — \code{ui\_dashboard\_system.c}}
\label{sec:system}

Tab System hiển thị trạng thái thiết bị, cập nhật mỗi 2\,s qua
\code{sys\_info\_timer\_cb} (\code{main.c:129} + \code{ui\_dashboard\_system.c:87}):
\begin{itemize}
\item \textbf{Heap}: \code{esp\_get\_free\_heap\_size} /
  \code{esp\_get\_minimum\_free\_heap\_size} (KB);
\item \textbf{Uptime}: \code{format\_uptime} từ \code{esp\_timer\_get\_time}
  (h:mm:ss);
\item \textbf{FW / IDF / build / flash}.
\end{itemize}

\noindent\textbf{Bảo mật token (R1):} \code{mask\_secret} (\code{ui\_dashboard\_system.c:53})
che chuỗi token; dòng \code{ui\_dashboard\_system.c:128} chỉ hiển thị
\code{coreiot\_token\_display()} đã được che — \textbf{không} bao giờ in token
đầy đủ lên màn hình.

\section{ESP-NOW status (LINKED)}
\label{sec:linked}

Hàm \code{ui\_dashboard\_set\_espnow\_status(linked)} (\code{ui\_dashboard.c:330})
đặt nhãn \code{"ESP-NOW: LINKED"} hoặc \code{"NO LINK"}.  Được gọi từ:
\begin{itemize}
\item \code{on\_espnow\_rx} — bất kỳ frame hợp lệ $\rightarrow$ LINKED;
\item \code{espnow\_link\_watchdog\_cb} — nếu hết link (xem $\S$\ref{sec:link-watchdog})
  $\rightarrow$ NO LINK + clear các slot quá \code{ESPNOW\_LINK\_TIMEOUT\_MS}.
\end{itemize}

\section{Tóm tắt UI}
\label{sec:ui-summary}

\textbf{Chuỗi cập nhật:} \code{on\_espnow\_rx/on\_coreiot\_data}
$\rightarrow$ \code{ui\_dashboard\_update\_sensor / clear\_sensor}
($\S$\ref{sec:main-c}) $\rightarrow$ \code{arc\_set\_zone}
$\rightarrow$ \code{evaluate\_hazard()} $\rightarrow$ banner OVERALL.
Mọi vòng lặp đều nằm trên LVGL task (qua \code{esp\_lv\_adapter\_lock()}),
đảm bảo thread-safe với MQTT/Wi-Fi.
